%!TEX root = ./main.tex

\section[Wrap-up]{Concluding Remarks}

% TALK MAPPING: slide 15, transcript 16:54--17:15, followed by the two questions the
% SIGCOMM audience actually asked (17:35--19:42). The Q&A frames are kept because both
% questions are the ones a student will ask too, and the speaker answered both.

% Teaching note (60 s): the authors' own summary, in their order. Do not add claims.
\begin{frame}{Concluding remarks}
  \vspace{0.8em}
  {\small\textbf{RANBooster: a fronthaul middlebox framework}\par}
  \vspace{0.2em}
  {\small\hspace{1.0em}$\cdot$~Build novel monitoring and control use cases\par}
  {\small\hspace{1.0em}$\cdot$~Accelerate innovation in the RAN space\par}

  \vspace{0.7em}
  {\small\textbf{Open source reference implementation}\par}

  \vspace{0.7em}
  \centering
  \begin{tcolorbox}[colback=white,colframe=softgreen,boxrule=0.8pt,arc=2pt,
    left=6pt,right=6pt,top=4pt,bottom=4pt,width=0.78\textwidth]
    \centering\small
    \href{https://github.com/microsoft/RANBooster}{\texttt{github.com/microsoft/RANBooster}}
  \end{tcolorbox}

  \vspace{0.6em}
  \takeaway{Four primitives on an interface every vendor already speaks\\
    were enough to add features the vendors never shipped.}
\end{frame}

% Teaching note (90 s): ask the class the question before revealing the answer. Both
% of these came from the room at SIGCOMM, and both are good instincts.
\begin{frame}[t]{Questions the SIGCOMM audience asked}
  \vspace{0.5em}
  \qa{The DU and the RUs are connected through a fronthaul switch. Where does
      this software actually run?}
     {In software on a server today. Offloading to SmartNICs or programmable
      switches is the obvious next step --- but packet \emph{caching} (A3) is the
      action that is not clearly implementable on a switch. Open problem.}

  \vspace{0.6em}
  \only<2->{%
    \qa{You said a middlebox on a latency-sensitive path might be a problem.
        Was it?}
       {No. The actions mostly rewrite header fields, which costs microseconds.
        That is what allowed scaling to eight or nine radios, and chaining one
        middlebox per floor underneath a building-wide DAS.}}
\end{frame}

% Teaching note (180 s): the discussion frame. Nothing here is in the authors' deck;
% these are the questions this course wants students to carry out of the paper.
\begin{frame}{What this paper is really arguing}
  \vspace{0.5em}
  \begin{columns}[T,onlytextwidth]
    \column{0.47\textwidth}
      {\small\textbf{\textcolor{ranblue}{The move}}\par}
      \vspace{0.3em}
      {\small
       \begin{itemize}\itemsep4pt
         \item The official extension point (the RIC) went unsupported.
         \item So use the interface vendors \emph{had} to implement.
         \item Middleboxes: the Internet's own answer to a frozen architecture.
       \end{itemize}\par}
    \column{0.47\textwidth}
      {\small\textbf{\textcolor{coreorange}{Worth arguing about}}\par}
      \vspace{0.3em}
      {\small
       \begin{itemize}\itemsep4pt
         \item Middleboxes made the Internet hard to evolve. Same trap here?
         \item Who operates a box that sees every user's raw waveform?
         \item Does this scale past one enterprise floor?
       \end{itemize}\par}
  \end{columns}

  \vspace{0.7em}
  \ask{Where else in this course have we seen a closed interface\\
    pried open by someone standing on the wire?}
\end{frame}
